\section{Evaluating Both Boundaries}\label{sec:reporting}
An intervention needs one record linking the attacker's remaining observations to the original task and allowed collateral effects. Table~\ref{tab:availability} shows why this joint record matters: among four query/workflow privacy settings from two sources, adaptation is documented in two and matched original-task utility in one, with no setting documenting both. The adapted settings share a preprint; the utility observation is WRA's published visibility comparison. These are descriptive counts of inspected configurations.

\begin{table*}[t]\centering\small
\caption{Evidence located in the retained records. Counts use eligible settings; the second column also gives distinct sources. ``Not located'' leaves evidence unresolved within the inspected scope. These counts describe availability, not study quality or checklist validity.}\label{tab:availability}
\setlength{\tabcolsep}{4pt}
\begin{tabularx}{\textwidth}{@{}P{.28\textwidth}P{.15\textwidth}P{.11\textwidth}Y@{}}\toprule
Question & Eligible settings / sources & Located & Supporting settings and consequence \\\midrule
Protected-condition attack adaptation & 4 / 2 & 2/4 & Communication graphs: transformed captures and live OHTTP. Both come from one preprint. \\\addlinespace[3pt]
Original-task utility under the same intervention & 4 / 2 & 1/4 & WRA visibility comparison. It does not price WRA's executed-decoy condition. \\\addlinespace[3pt]
Adaptation and original-task utility in one setting & 4 / 2 & 0/4 & Adaptation and utility are located in different records; transport cost alone does not fill the task-outcome field. \\\addlinespace[3pt]
Original-task utility for executed cover & 1 / 1 & 0/1 & WRA's five research decoys report recovery and extra research; matched original-task utility is not located. \\\bottomrule
\end{tabularx}
\end{table*}


\subsection{A minimal evaluation procedure}
The following procedure connects established evaluation guidance~\cite{arp2022dos,juarez2014critical,mathews2023critical} to the information and action boundaries. It is a proposed reporting procedure, illustrated by the controlled cover case. The six observation-contract fields in Section~3 specify an inference setting; the eight checklist items organize the broader study report, including its outcomes, uncertainty and provenance.

\begin{enumerate}
\item \textbf{Declare the target and permitted observations.} Name the private fact, visible model/tool/recipient interfaces, access rights and any supplied execution windows. Track where tool results are forwarded.
\item \textbf{Declare the allowed execution effects.} Specify the original goal, other objects that may change, resource budgets and what cover may return to the planner. Use business invariants that permit legitimate concurrent updates; exact state equality fits only a correspondingly constrained fixture.
\item \textbf{Pair protection, attacks and outcomes.} Evaluate the baseline and intervention on matched tasks. Train permitted attackers on protected executions and include attacks that exploit the cover schedule. Record goal completion, completed tasks violating the state contract, and bytes, delay and resource use on those same runs.
\item \textbf{Make the decision against stated tolerances.} Report attack-specific inference and uncertainty together with task, state and resource tolerances fixed before evaluation. A completed task with a state change outside the declared contract fails its utility requirement. Successful inference on another permitted view defeats a claim of containment across those views. An unsuccessful tested attack supplies no general secrecy guarantee.
\end{enumerate}

\paragraph{Worked application}
The cover study binds requested-action recognition to a tool-metadata observer and checks both target completion and the whole synthetic fixture. Its adverse live-write arm separates target completion from collateral effects, while the other policies record residual recognition and cost under the same observer. Table~\ref{tab:cover-main} keeps the stronger position rule visible. Supplement S8 derives the fixture contract and fills each procedure step. The application demonstrates joint reporting; model/timing views and deployment acceptance thresholds require additional evidence.

\subsection{Corpus availability and reusable forms}
Three identity-obfuscation settings, two adversarial-Agent blocking settings and thirty-five without a protection comparison have different eligibility from the four privacy settings. The access map separately assigns twenty-two transit, ten recipient, eleven active and one downstream preference setting. Supplement S4.1 reports outcome/failure encoding sensitivity with each encoding's own denominator. Supplement S5 provides the eight-field checklist, its correspondence to established guidance and filled WRA, OHTTP and ANT examples. Each form records evidence, a missing item or inapplicability.
